\chapter{The Pipeline in Seven Stages}
\label{ch:pipeline}

\section{Method}

The chapters of Part~II describe the instruments one at a time. This chapter shows
them accumulating on a single board.

The subject is Whitegape, the destroy board of \cref{ch:dtm}. Its finish was
truncated to a prefix of its keys seven times, and each truncation was driven
through the studio under a slug of its own, exported, and rendered from the same
camera. Nothing between the pictures was drawn by hand, and nothing about the plan
changed: the only difference between one panel and the next is that one more key
of the finish was present.

\begin{figure}[htbp]
\centering
\begin{subfigure}{0.48\textwidth}\includegraphics[width=\linewidth]{stages/01-ground.png}
  \caption{Ground. The plan compiled and nothing else.}\end{subfigure}\hfill
\begin{subfigure}{0.48\textwidth}\includegraphics[width=\linewidth]{stages/02-relief.png}
  \caption{Relief. Four marks and their pushes carve the fell.}\end{subfigure}

\medskip
\begin{subfigure}{0.48\textwidth}\includegraphics[width=\linewidth]{stages/03-paint.png}
  \caption{Paint. A theme per bucket, banded by the ground's own angle.}\end{subfigure}\hfill
\begin{subfigure}{0.48\textwidth}\includegraphics[width=\linewidth]{stages/04-shapes.png}
  \caption{Shapes. Nine authored: the pit, the dock, two flights, two ramps, a wall,
  two kerbs.}\end{subfigure}

\medskip
\begin{subfigure}{0.48\textwidth}\includegraphics[width=\linewidth]{stages/05-layers.png}
  \caption{Layers. One made layer, the kiln, at base 18.}\end{subfigure}\hfill
\begin{subfigure}{0.48\textwidth}\includegraphics[width=\linewidth]{stages/06-rooms.png}
  \caption{Rooms. The spawn shell, in place of the built-in box.}\end{subfigure}

\medskip
\begin{subfigure}{0.48\textwidth}\includegraphics[width=\linewidth]{stages/07-dressed.png}
  \caption{Dressed. Ten styles and 27 props.}\end{subfigure}
\caption{One board, seven exports. Each panel adds one key of the finish to the one
before it.}
\label{fig:stages}
\end{figure}

\section{What Each Stage Adds}

\Cref{tab:stages} is the same sequence as numbers: what the finish carried, how many
shapes the stored layout held, how many distinct surface blocks the built world's
theme census counted, and what the store and the export answered with.

\begin{table}[htbp]
\caption{The seven stages, measured. The key added at each stage is bold; the rest
is inherited.}
\label{tab:stages}
\scriptsize
%% Generated from seven runs of tools/drive.py against the running pgm-studio API
%% (http://localhost:7894/api) --- do not edit by hand.
%% One row per pipeline stage of specs/opus5-whitegape, each driven under its own slug
%% (stage-01-ground ... stage-07-dressed) with the finish truncated to the keys named.
%% The key added at each stage is bold; the rest is what the stage inherited.
%% Shapes: shapes in the stored layout, over every layer. Blocks: distinct surface blocks
%% the built world's theme census counted (GET /map/{slug}/themes/census).
%% Findings: every rule id raised on a 2xx, grouped by the route that raised it. Every stage also
%% drew G8 (dead-share 0.198 outside the authored band) from POST /plan/evaluate, which
%% judges the plan --- identical in all seven runs --- and not the finish.
%% A key the spec does not state is carried by no stage, so it appears in no row.

\begin{tabularx}{\linewidth}{@{}l >{\raggedright\arraybackslash}X r r >{\raggedright\arraybackslash}p{27mm}@{}}
\toprule
\textbf{Stage} & \textbf{What the finish carried} & \textbf{Shapes} & \textbf{Blocks} & \textbf{Findings} \\
\midrule
\id{01-ground} & \textbf{\id{created}}, \textbf{\id{authors}}\newline\textit{the plan compiled and nothing else: bedrock and stone at the surface the plan states} & 6 & 4 & \id{SK8}\ \textit{\scriptsize(the store)} \\
\addlinespace[2pt]
\id{02-relief} & \id{created}, \id{authors}, \textbf{\id{relief}}\newline\textit{5 marks and 3 pushes carve the fell} & 6 & 4 & \id{RL2}, \id{RL3}\ \textit{\scriptsize(relief read)} \\
\addlinespace[2pt]
\id{03-paint} & \id{created}, \id{authors}, \id{relief}, \textbf{\id{themes}}, \textbf{\id{mapTheme}}, \textbf{\id{biome}}\newline\textit{5 themes, the map default and a solid biome} & 6 & 8 & \id{RL2}, \id{RL3}\ \textit{\scriptsize(relief read)} \\
\addlinespace[2pt]
\id{04-shapes} & \id{created}, \id{authors}, \id{relief}, \id{themes}, \id{mapTheme}, \id{biome}, \textbf{\id{addShapes}}, \textbf{\id{editShapes}}, \textbf{\id{bendShapes}}\newline\textit{15 authored shapes, plus the coast edited and bent} & 21 & 10 & \id{RL2}, \id{RL3}\ \textit{\scriptsize(relief read)} \\
\addlinespace[2pt]
\id{05-layers} & \id{created}, \id{authors}, \id{relief}, \id{themes}, \id{mapTheme}, \id{biome}, \id{addShapes}, \id{editShapes}, \id{bendShapes}, \textbf{\id{addLayers}}\newline\textit{2 made layers: kiln-base, kiln-stack} & 29 & 10 & \id{RL2}, \id{RL3}\ \textit{\scriptsize(relief read)} \\
\addlinespace[2pt]
\id{06-rooms} & \id{created}, \id{authors}, \id{relief}, \id{themes}, \id{mapTheme}, \id{biome}, \id{addShapes}, \id{editShapes}, \id{bendShapes}, \id{addLayers}, \textbf{\id{roomStyles}}\newline\textit{the spawn shell, in place of the built-in box} & 29 & 10 & \id{RL2}, \id{RL3}\ \textit{\scriptsize(relief read)} \\
\addlinespace[2pt]
\id{07-dressed} & \id{created}, \id{authors}, \id{relief}, \id{themes}, \id{mapTheme}, \id{biome}, \id{addShapes}, \id{editShapes}, \id{bendShapes}, \id{addLayers}, \id{roomStyles}, \textbf{\id{dressing}}\newline\textit{10 styles and 19 props} & 29 & 17 & \id{RL2}, \id{RL3}\ \textit{\scriptsize(relief read)} \\
\addlinespace[2pt]
\bottomrule
\end{tabularx}

\end{table}

\section{Four Observations}

\subsection{The First Panel Is What a Plan Actually Buys}

Stage one is the plan compiled with an empty finish: two flat grey slabs, a gap
between them, and the objective markers. It is worth looking at for a moment,
because it is exactly what \cref{ch:division} claims a composed or authored plan
provides. The arrangement is entirely present. The place is entirely absent.

The spawn and wool shells are already standing, which is easy to misread as the
finish having leaked into the stage. It has not. Those structures are stamped from
the \emph{plan's} own pieces rather than from the finish's room styles, which is
why stage six, where the room style arrives, changes their material and not their
existence.

\subsection{The Block Count Is the Whole Story in One Column}

The census of distinct surface blocks runs 4, 4, 8, 10, 10, 10, 11 across the seven
stages. Two things fall out of that sequence.

Relief adds nothing to it. Stage two reshapes the entire board and leaves the block
count exactly where stage one had it, because relief decides where ground goes and
never what it is made of. The two concerns are orthogonal in the model and the
measurement shows it.

Paint doubles it, from four to eight, in one key. Then authored shapes add two more,
and everything after that adds one between them. By far the largest single change in
what the board is \emph{made of} is the theme registry, and by far the largest
single change in what the board \emph{is} comes later.

\subsection{Most Stages Are Silent}

Only two of the seven answered with any finding at all. Stage one drew \id{SK8}.
Stage seven drew six rules at once, every one of them from the dressing family:
\id{DR-CLAIM}, \id{DR-CUT}, \id{DR-KEEP}, \id{DR-PASS}, \id{DR-ROAD} and
\id{DR-STEEP}.

That distribution is not an accident of this board. Dressing is the stage that
places things on ground it did not make, so it is the stage with the most to
complain about: a prop claiming a cell another prop claimed, a prop cut in two by
an obstruction, a prop on ground too steep to stand on, a road crossing something
it should not. Everything before it is authored against an empty board and has
nothing to collide with.

\subsection{The Plan's Own Verdict Never Moves}

Every one of the seven runs drew the same plan-level complaint: a dead-share of
0.198, outside its authored band. It is identical in all seven because it judges
the \emph{plan}, and the plan did not change.

This is a useful demonstration of a distinction \cref{ch:refusals} makes in the
abstract. A plan-tier rule reads pieces and build zones. It cannot see relief, paint,
layers, rooms or props, so no amount of finishing will move it. An author trying to
answer a plan-tier complaint by improving the finish is working at the wrong level,
and the seven identical readings are what that looks like from outside.
